% Generated draft: requires mathematical and source review.
\documentclass[11pt]{article}
\usepackage[margin=25mm]{geometry}
\usepackage{fontspec}
\setmainfont{Latin Modern Roman}
\usepackage{amsmath,amssymb,mathtools,unicode-math}
\setmathfont{Latin Modern Math}
\usepackage{enumitem,needspace,xurl,hyperref}
\setlength{\parindent}{0pt}
\setlength{\parskip}{5pt}
\setlength{\emergencystretch}{3em}
\newcommand{\N}{\mathbb{N}}
\newcommand{\Z}{\mathbb{Z}}
\newcommand{\Q}{\mathbb{Q}}
\newcommand{\R}{\mathbb{R}}
\newcommand{\C}{\mathbb{C}}
\newcommand{\F}{\mathbb{F}}
\newcommand{\A}{\mathbb{A}}
\newcommand{\PP}{\mathbb{P}}
\newcommand{\E}{\mathbb{E}}
\newcommand{\eps}{\varepsilon}
\newcommand{\dd}{\,\mathrm{d}}
\begin{document}
\section*{UnsolvedMath ID: 30004122}
\textbf{Source problem code:} catalogue entry
\subsection{Definitions and mathematical statement}
For a finite undirected simple graph \(F\), define
\[
\operatorname{ex}(n,F)=\max\{e(G): |V(G)|=n,\ G\text{ is a finite undirected simple graph containing no subgraph isomorphic to }F\},
\]
whenever the family over which the maximum is taken is nonempty. Containment here means ordinary, not necessarily induced, subgraph containment. These are the conventional graph and extremal-number interpretations used in this draft; the supplied statement does not explicitly specify them.

For an even integer \(k\geq2\), construct \(H^{k-1}\) by replacing each edge \(uv\) of \(H\) with a path from \(u\) to \(v\) having exactly \(k\) edges and \(k-1\) new internal vertices. All new internal vertices are distinct across replacement paths and lie outside \(V(H)\); isolated vertices of \(H\) are retained.

The supplied \emph{Conlon--Lee subdivision conjecture} asks whether, for every even integer \(k\geq2\) and every graph \(H\), there exists \(\delta=\delta(k,H)>0\) such that
\[
\operatorname{ex}(n,H^{k-1})=O\!\left(n^{1+1/k-\delta}\right)
\qquad(n\to\infty).
\]
Explicitly, for each fixed pair \((k,H)\), the requested conclusion is that there exist real constants \(\delta>0\), \(C>0\), and an integer \(n_0\geq1\) such that, for every integer \(n\geq n_0\),
\[
\operatorname{ex}(n,H^{k-1})\leq Cn^{1+1/k-\delta}.
\]
The constants may depend on \(k\) and \(H\), but not on \(n\).

\emph{Unresolved boundary qualification.} The source quantifies over every graph \(H\), without excluding edgeless graphs. Under ordinary subgraph containment, if \(H\) is edgeless and \(n\geq |V(H)|\), every \(n\)-vertex graph contains \(H^{k-1}=H\). The defining family for the maximum is then empty. Thus the universal wording requires either an explicit convention for this case or a hypothesis that \(H\) has an edge; neither is supplied. For finite simple \(H\) with at least one edge, the displayed formulation is well-defined.

The supplied dated curation labels the record partially solved and reports the bipartite case as proved and the general conjecture as open. It also describes a separate \(L\)-subdivision problem as completely resolved by Corollary 13. That separate problem is not formulated in the supplied record. These are supplied source assessments, not independently verified mathematical or literature conclusions.
\subsection{Short English statement}
For each fixed even integer \(k\geq2\) and graph \(H\), the conjecture asks whether forbidding the subdivision obtained by replacing each edge of \(H\) with a path of \(k\) edges forces an extremal bound with exponent strictly below \(1+1/k\). The improvement may depend on \(k\) and \(H\). The graph conventions and edgeless boundary case require confirmation; the supplied partial-resolution assessment remains unverified.
\subsection{Sources}
\begin{itemize}
\item Local registry: \texttt{data/unsolvedmath/problems.json}, with \texttt{catalog/problems.json} as the additional-record registry; supplied permanent UnsolvedMath ID \(30004122\), exact source code \texttt{catalogue entry}, title \emph{Conlon--Lee Conjecture for Graph Subdivisions}. These registry locations are cited as project provenance and were not inspected for this draft.
\item Supplied report citation: \emph{Combinatorics, Probability and Computing} (2020), Oberwolfach Reports, pp.\ 1113--1189, DOI \texttt{10.4171/owr/2019/19}: \url{https://doi.org/10.4171/owr/2019/19}. The record associates the problem with 2019 and attributes the conjecture to Conlon--Lee; no separate proposer is specified.
\item Supplied literature source: Oberwolfach Report 19/2019, Theorem 12 and Corollary 13: \url{https://doi.org/10.4171/owr/2019/19}.
\item Supplied literature source: Conlon--Lee, \emph{On the extremal number of subdivisions}: \url{https://doi.org/10.1093/imrn/rnz088}.
\item Supplied literature source: Conlon--Lee--Janzer, \emph{More on the extremal number of subdivisions}: \url{https://doi.org/10.1007/s00493-020-4202-1}.
\item Literature assessment dated exactly \texttt{2026-08-22}: the supplied curation labels the record partially solved, reports the conjecture proved for bipartite graphs and open in general, and reports a separate displayed \(L\)-subdivision problem resolved by Corollary 13. This is supplied dated curation, not independently verified; its time-sensitive status and cited claims require independent checking.
\item The supplied record labels its statement \texttt{corrected\_verified} and says that the official report was checked directly and definitions restored. This records the curator's verification claim; no source visit or independent verification was performed for this draft.
\end{itemize}
\end{document}
